\chapter*{Conclusion Générale}
\addcontentsline{toc}{chapter}{Conclusion Générale}
\markboth{Conclusion Générale}{Conclusion Générale}

\begin{resultatbox}[Bilan du projet]
Ce projet a produit une plateforme MLOps complète, industrialisée et documentée
pour la détection précoce du risque d'hospitalisation. L'ensemble du cycle de
vie d'un modèle de machine learning --- de la donnée brute au service de
production surveillé --- a été implémenté et validé.
\end{resultatbox}

\vspace{1em}

\section*{Contributions majeures}

Ce projet se distingue d'un projet de machine learning classique par sept
contributions concrètes~:

\begin{enumerate}

\item \textbf{Split temporel strict 2008 / 2009 / 2010.}
  Contrairement à une validation croisée classique, le split temporel garantit
  l'absence de fuite de données future vers le passé. L'AUC de 0.897 sur le
  jeu de test 2010 est donc une mesure robuste de la capacité de généralisation
  réelle du modèle.

\item \textbf{Construction anti-leakage de la variable cible.}
  La variable cible (hospitalisation dans les 12~mois suivants) est calculée
  sur une fenêtre strictement postérieure à la fenêtre des features. Ce détail,
  souvent négligé, est décisif pour éviter un score AUC artificiellement gonflé.

\item \textbf{Optimisation du seuil de décision.}
  Le seuil $\tau^* = 0.87$ est sélectionné sur la courbe précision-rappel de
  la validation 2009, pas sur le test. Cette démarche est documentée et
  reproductible via MLflow.

\item \textbf{Monitoring avec rollback automatique réellement testé.}
  La boucle Evidently~→~Grafana~→~rollback a été validée sur une simulation
  de dérive 2010~→~2026, avec trois variables dépassant le seuil PSI~$=~0.20$.

\item \textbf{Analyse de fairness par groupe racial.}
  Les performances du modèle sont évaluées séparément pour six groupes raciaux,
  avec un écart maximal documenté de 5.5\,\% en rappel.

\item \textbf{ROI calculé et justifié financièrement.}
  Le ROI net de \$2.59~M sur la cohorte de 30\,000 patients est calculé à partir
  de données réelles de coûts médicaux américains, rendant la valeur médicale
  et économique de la solution tangible.

\item \textbf{Simulation de dérive temporelle 2010~→~2026.}
  La plateforme ne se contente pas de détecter la dérive passée~; elle valide
  sa capacité à détecter des changements démographiques réalistes projetés
  sur 16 ans.

\end{enumerate}

\section*{Ce qui distingue cette approche}

Là où beaucoup de projets de machine learning médical s'arrêtent à la publication
d'un AUC, cette plateforme répond à la question qui compte vraiment~:
\textit{le modèle fonctionne-t-il encore dans un an~?} La réponse est oui,
grâce à une architecture de monitoring qui ne dépend d'aucune intervention
humaine pour détecter et corriger la dérive.

L'analyse de fairness et le calcul ROI transforment une métrique technique
en un argument actionnable pour les décideurs de santé publique~: investir dans
la prévention guidée par l'IA génère un retour financier mesurable tout en
réduisant les inégalités de soin.

\section*{Perspectives futures}

Ce projet ouvre plusieurs axes d'amélioration~:

\begin{itemize}
  \item \textbf{Passage à l'échelle~:} déploiement sur les 20~samples
        (2.3~millions de patients) avec Apache Spark et Kubernetes.
  \item \textbf{Réduction de l'écart de fairness~:} contraintes de fairness
        intégrées dans la fonction objectif d'Optuna.
  \item \textbf{Modèles de prédiction à horizon variable~:} prédire le risque
        à 30, 90 et 180~jours en parallèle.
  \item \textbf{Intégration au DMP~:} connexion aux Dossiers Médicaux Partagés
        pour des features plus riches (ordonnances réelles, résultats biologiques).
  \item \textbf{Explicabilité locale à la demande~:} endpoint \texttt{/explain}
        qui retourne les valeurs SHAP individuelles pour chaque prédiction.
\end{itemize}

\vspace{1em}
\noindent Ce projet démontre que le MLOps n'est pas une couche optionnelle~:
c'est ce qui transforme un modèle expérimental en un outil de soin fiable.
